Hasty is a pioneer and leader in synthetic biology. His early career focused on tackling central impediments to the creation of an engineering discipline. He created genetic clocks to develop the underlying equations that drive circuit design, and used them to explore the coupling of gene circuits with their host genome. Faced with the inevitable noise of intracelluar gene regulation, he was the first to demonstrate how bacterial colonies can be engineered to function deterministically through the use of intercellular synchronization. He introduced the concept of “synergistic synchronization,” whereby two synchronization mechanisms couple colonies of bacteria at centimeter length scales. He used this concept to develop inexpensive biosensors that don't require complex optics. He has engineered periodic lysis of a bacterial colony such that population levels oscillate within a tumor microenvironment and release an encoded therapeutic. This design directly addresses the problem of systemic inflammatory response with programmed population control; since the colony is pruned after each oscillatory lysis event, the design mitigates an undesirable host response.      

Jeff Hasty is a pioneer and leader in synthetic biology. His early career focused on tackling central impediments to the creation of an engineering discipline. He created genetic clocks to develop the underlying equations that drive circuit design, and used them to explore the coupling of gene circuits with their host genome. Faced with the inevitable noise of intracelluar gene regulation, he was the first to demonstrate how bacterial colonies can be engineered to function deterministically through the use of intercellular synchronization. He introduced the concept of “synergistic synchronization,” whereby two synchronization mechanisms couple colonies of bacteria at centimeter length scales. He used this concept to develop inexpensive biosensors that don't require complex optics. He has engineered periodic lysis of a bacterial colony such that population levels oscillate within a tumor microenvironment and release an encoded therapeutic. This design directly addresses the problem of systemic inflammatory response with programmed population control; since the colony is pruned after each oscillatory lysis event, the design mitigates an undesirable host response. He has engineered a naturally competent bacterium to take up and read mammalian DNA in the environment, and shown that this mechanism can be used to detect single-base mutations associated with colorectal cancer.  He has shown how his genetic clock designs can be used to engineer longevity in single cells. Cells and DNA constructs from his lab are used by scientists and engineers around the world in efforts to build upon his development of logical programming at the colony level. 

\cventry{1992--1997}
\
Ph.D. in Statistical Physics from the Georgia Institute of Technology (1992-97); Postdoc in Bioengineering at Boston University (1998-2002); Assistant Professor of Bioengineering, UC San Diego (2002-06); Associate Professor of Bioengineering, UC San Diego (2006-10); Associate Professor of Biology, UC San Diego (2009-2010); Director, Synthetic Biology Institute, UC San Diego (2009-present); Professor of Bioengineering and of Biology, UC San Diego (2010-present).  Research Foci: Systems and Synthetic Biology:  (i) Development of a theoretical understanding of living systems, (ii) Transitioning synthetic biology towards diagnostics and therapeutics, (iii) Systems engineering of water treatment, (iv) Development of biosensors for continuous monitoring of industrial processes, (v) transforming eduction for the next generation of biologists.  Finalist, Howard Hughes Medical Institute Investigator Program; Fellow, American Institute for Medical and Biological Engineering; Teacher of the Year, UCSD Bioengineering; Alfred P. Sloan Research Fellow; NSF Career Award.  Founder and lead organizer of Winter q-BIO Meeting (2012-present); Associate Editor, ACS Synthetic Biology (2011-17); Editorial Boards: Biophysical Journal (2010-12), Systems and Synthetic Biology (2006-12), Chaos (2007-15), IEE Systems Biology (2004-06).  Founder:  GenCirq Inc. and Quantitative BioSciences Inc.  


\item Founder and Lead Organizer -- Annual Winter Q-BIO Meeting (2012-present). 
\item Associate Editor, {\em ACS Synthetic Biology} (2011-2017)
\item Organizer (with Ron Weiss), Synthetic Biology Workshop.  Mathematical Biosciences Institute, Ohio State University, Columbus, OH (January, 2010).  
\item Editorial Board, {\em Biophysical Journal} (2010-2012). 
\item Editorial Board, {\em Systems and Synthetic Biology} (2006-2012).
\item Editorial Board, {\em Chaos} (2007-2015).
\item Organizer - Minisymposium: Dynamics in Gene Expression, Society for Industrial
and Applied Mathematics (SIAM) 2nd Life Science Conference, Raleigh,
NC (2006).
\item Workshop Organizer, Columbus, OH. Genomics/proteomics/bioinformatics
workshop at the NSF Mathematical Biosciences Institute (11/08/2004).
\item Computational Bioengineering Track Organizer~(9 sessions), Biomedical Engineering Society (BMES) annual meeting, Philadelphia, PA (10/14/2004).
\item Organizer.  Minisymposium on dynamics of gene regulation.
Society for Industrial and Applied Mathematics (SIAM) 2nd Life
Science Conference, Portland, OR (7/11/04).
\item Discussion Leader: Modeling Transcriptional Control in Gene Regulatory Networks,
Gordon Research Conference: Theoretical Biology \& Biomathematics
(6/9/04).
\item Editorial Board, {\em IEE Systems Biology} (2004-2006).
\item Scientific Advisory Committee, CIIT Centers for Health
Research, Research Triangle Park, NC (2004-2006).
\item Chair -- Session: Noise in genetic networks. American Physical
Society March Meeting, Austin, TX (2003).
\item NSF Delegation Member. Served as one of three junior scientists
on a US team formed to participate in the Advanced Studies
Institute in Canberra Australia (2002).
\item Organizer - Minisymposium: Gene regulatory networks, Society for Industrial
and Applied Mathematics (SIAM) 1st Life Science Conference,
Boston, MA (2002).
\item Organizer - Minisymposium: Dynamics of gene networks, Society for Industrial and
Applied Mathematics (SIAM) 6th Conference on Applications of
Dynamical Systems, Snowbird, UT (2001).
\item Chair - Session: Designer gene networks: Towards fundamental cellular control,
Society for Industrial and Applied Mathematics (SIAM) 6th
Conference on Applications of Dynamical Systems, Snowbird, UT
(2001).
\end{itemize}




Finalist, Howard Hughes Medical Institute Investigator Program (2015). 

\item Fellow, American Institute for Medical and Biological Engineering (2014). 

\item Teacher of the Year, UCSD Department of Bioengineering (2010).

\item Invited Speaker -- Nobel Symposium on Systems Biology~(2009).

\item Alfred P. Sloan Research Fellow~(2003--05).

\item NSF Career Award (2003--08).

 

\section{Research Focus: Systems and Synthetic Biology}
\begin{tabular}{lll} 
\href{http://qbio.ucsd.edu}{Development of a theoretical understanding of living systems}\\
\href{http://gencirq.com}{Transitioning synthetic biology towards diagnostics and therapeutics for cancer}\\
\href{https://www.qbisci.com/fiscalini-dairy}{Systems engineering of water treatment for the generation of feedstock, energy and clean water}\\
\href{https://www.qbisci.com/sensors/}{Development of biosensors for continuous monitoring of industrial processes}\\
\href{https://qbio.ucsd.edu}{Transforming education for the next generation of biologists}\\
\href{http://biocircuits.sdsc.edu/outreach/}{Integration of new technologies into K6 science education}\\
\end{tabular}\\
    


 % arguments 3 to 6 can be left empty
%\cventry{1988--1991}{M.S. Physics}{San Francisco State University}{Atlanta}{GA}{Advisor: Kurt Wiesenfeld}

%\cventry{year--year}{Degree}{Institution}{City}{\textit{Grade}}{Description}

%\section{Master thesis}
%\cvitem{title}{\emph{Title}}
%\cvitem{supervisors}{Supervisors}
%\cvitem{description}{Short thesis abstract}

\section{Academic Appointments}
 \cventry{2010--present}{Professor}{Department of Bioengineering}{UC San Diego}{}{~}
\vspace{-.2in}\cventry{2010--present}{Professor}{Molecular Biology Section}{UC San Diego}{}{~}
\vspace{-.2in} \cventry{2009--present}{Director}{Synthetic Biology Institute}{UC San Diego}{}{~}
\vspace{-.2in}\cventry{2009--2010}{Associate Professor}{Molecular Biology Section}{UC San Diego}{}{~}
\vspace{-.2in}\cventry{2006--2010}{Associate Professor}{Department of Bioengineering}{UC San Diego}{}{~}
\vspace{-.2in}\cventry{2002--2006}{Assistant Professor}{Department of Bioengineering}{UC San Diego}{}{~}
\vspace{-.2in}\cventry{2001--2002}{Research Assistant Professor}{Department of Bioengineering}{Boston University}{}{~}
\vspace{-.2in}\cventry{1998--2001}{Postdoc with Jim Collins}{Department of Bioengineering}{Boston University}{}{~}

